\nottoggle{beamer}{
  % enumerate with Roman numerals
  \newlist{renum}{enumerate}{1}
  \setlist[renum,1]{label=(\roman*)}
  % enumerate alphabetically
  \newlist{aenum}{enumerate}{1}
  \setlist[aenum,1]{label=(\alph*)}
}{}

% Introducing a 'dummyCounter' to be used whenever no explicit counter is given. 
% We initialize the dummyCounter to a large number as to not interfere
% with explicit counters (which should never realistically reach such large numbers).
\newcounter{dummyCounter} \setcounter{dummyCounter}{99999}
\newcommand{\mylabel}{dummyLabel} % will get overwritten

% within chapter / section counter (cc = chapter counter)
\iftoggle{book}{
  \newcounter{cc}[chapter] % cc shall be reset at the start of a new chapter
}{
  \newcounter{cc}[section] % cc shall be reset at the start of a new section
}

\setcounter{cc}{0}

\iftoggle{book}{
  \renewcommand{\thecc}{\arabic{chapter}.\arabic{cc}}
}{
  \renewcommand{\thecc}{\arabic{section}.\arabic{cc}}
}

% generic tcolorbox template for aufgabe, myantwort, myexample and myachtung
\newtcolorbox{mybox}[1][]{
  enhanced jigsaw,
  lines before break=9,
  breakable,
  attach boxed title to top text left={yshift=-3mm,yshifttext=-3mm},
  #1
}

\newaliascnt{ccaufgabe}{cc}
\aliascntresetthe{ccaufgabe}
\newenvironment{aufgabe}[2][]
{
  \refstepcounter{ccaufgabe}
  \begin{tcolorbox}[enhanced jigsaw,
    colback=RoyalBlue!10,
    colframe=RoyalBlue,
    colbacktitle=RoyalBlue,
    breakable,
    lines before break=9,
    title=\faEdit{} Aufgabe~\theccaufgabe]
    \ifblank{#2}{
      % #2 is blank
      \renewcommand{\mylabel}{\theccaufgabe}
      % \refstepcounter{cc}
    }{
      % #2 is an actual user-defined label
      \renewcommand{\mylabel}{#2}
    }
    \label{\mylabel}
    } %%
    {
  \end{tcolorbox}
  }

\newenvironment{myantwort}[1][]
{
\renewcommand{\AnswerHeader}{}
\begin{Answer}[ref=\theccaufgabe]
  \begin{mybox}[
    colback=ForestGreen!10,
    opacityback=.5,
    coltitle=white,
    colframe=ForestGreen,
    colbacktitle=ForestGreen,
    opacityback=1,
    title=\faCheckSquareO{} Lösungsvorschlag zu Aufgabe \theccaufgabe]
    } %%
    {
  \end{mybox}
\end{Answer}}

% TODO: add one more argument to give label (just like in aufgabe)
\newaliascnt{ccabgabe}{cc}
\aliascntresetthe{ccabgabe}
\newenvironment{myabgabe}[1][]
{
  \refstepcounter{ccabgabe}
  \begin{tcolorbox}[enhanced jigsaw,
    colback=Purple!10,
    colframe=Purple,
    colbacktitle=Purple,
    breakable,
    lines before break=9,
    title=\faArrowUp{} Aufgabe (Abgabe)~\theccabgabe]
    } %%
    {
  \end{tcolorbox}
}

% TODO: add one more argument to give label (just like in aufgabe)
\newaliascnt{ccchallenge}{cc}
\aliascntresetthe{ccchallenge}
\newenvironment{mychallenge}
{
  \refstepcounter{ccchallenge}
  \begin{tcolorbox}[enhanced jigsaw,
      colback=Gold!10,
      colframe=Gold!50!black,
      colbacktitle=Gold!50!black,
      breakable,
      lines before break=9,
      title=\faTrophy{} Challenge-Aufgabe~\thecc]
    } %%
    {
  \end{tcolorbox}
}

\newenvironment{myachtung}
{
  \begin{mybox}[
      colback=Red!10, % Background color
      coltitle=white,
      colframe=Red!50!black,
      colbacktitle=Red!50!black,
      title=\faWarning{} Achtung!]
    } %%
    {
  \end{mybox}
}

\newaliascnt{ccexample}{cc}
\aliascntresetthe{ccexample}
\ExplSyntaxOn
\NewDocumentEnvironment{myexample}{o}
{%
  \par
  \addvspace{.15in}%
  \refstepcounter{ccexample}%
  \IfNoValueTF{#1} {
    \noindent\textbf{Beispiel~\thecc.} \rmfamily
  }
  {
    \noindent\textbf{Beispiel~\thecc.} \label{ex:#1} \rmfamily
  }
}{%
  \par\addvspace{.15in}%
}
\ExplSyntaxOff

\newaliascnt{ccbemerkung}{cc}
\aliascntresetthe{ccbemerkung}
\ExplSyntaxOn
\NewDocumentEnvironment{mybemerkung}{ o }
{%
  \par
  \addvspace{.15in}%
  \refstepcounter{ccbemerkung}%
  \IfNoValueTF{#1} {
    \noindent\textbf{Bemerkung(en)~\thecc.} \rmfamily
  }
  {
    \noindent\textbf{Bemerkung(en)~\thecc.} \label{bemerkung:#1} \rmfamily
  }
}{%
  \par\addvspace{.15in}%
}
\ExplSyntaxOff

\nottoggle{beamer}{
  %% naming for clever refs in custom environments %%
  \crefname{ccaufgabe}{Aufgabe}{Aufgaben}
  \Crefname{ccaufgabe}{Aufgabe}{Aufgaben}
  \crefname{ccabgabe}{Abgabe}{Abgaben}
  \Crefname{ccabgabe}{Abgabe}{Abgaben}
  \crefname{ccchallenge}{Challenge}{Challenges}
  \Crefname{ccchallenge}{Challenge}{Challenges}
  \crefname{ccexample}{Beispiel}{Beispiele}
  \Crefname{ccexample}{Beispiel}{Beispiele}
  \crefname{ccexample}{Bemerkung}{Bemerkungen}
  \Crefname{ccexample}{Bemerkung}{Bemerkungen}
}{}
